\section{Crop classification from history and seasonal observations}
\label{sec:classification}

\subsection{What seasonal observations add}

We tested whether current-season NDVI and soil moisture improve crop
classification beyond the information in preceding crop labels.
The target was the current year's CDL class: corn (code~1), soybean (5),
winter wheat (24), or Other. Other combines the remaining selected CDL codes
(0--61), including code~0. The experiment used concurrent observations
through late season; forecasting before planting would require an earlier
observation cutoff and a separate evaluation.

The classification mask retained cells with selected CDL codes (0--61) in
at least three years during 2013--2022, separately from the rotation
analysis's corn--soybean eligibility rule. Each sample represented a grid
cell in one year. The full model used 38
features: 19 from CDL history, 15 from NDVI, and four from SMAP. CDL features
summarized up to ten preceding years through five annual lags, corn--soybean
transition counts, crop frequencies, time since each crop, sequence metrics,
and $3\times3$ neighborhood summaries. NDVI features described the current
season's level, amplitude, peak and green-up timing, rate of green-up, and
early-, middle-, and late-season means. Four additional NDVI summaries
described historical peak magnitude and timing. These features were calculated
directly from the weekly observations. SMAP supplied spring and growing-season
means plus fractions of weeks above an 80th-percentile or below a
20th-percentile reference.

We fitted a LightGBM classifier \citep{ke2017lightgbm} with a maximum of 500
boosting rounds, learning rate 0.05, maximum tree depth seven, and 63 leaves.
Training used 4.5 million samples from 2013--2021; the 2013--2014 rows retained
missing SMAP inputs because that record begins in 2015. The 2022 validation set
controlled early stopping after 50 rounds without improvement and supported
the feature-group comparison. The full model was then evaluated on 2023.
Sampling selected 125,000 observations from each class per year, giving
500,000 observations in each held-out year. Accuracy therefore measures
agreement with CDL on a balanced sample. Estimating regional map accuracy
would require area-weighted evaluation. Macro F1 averages the four
class-specific F1 scores with equal weight.

\begin{table}[htbp]
\centering
\small
\begin{tabular}{lrrr}
\hline
Features & Number & Accuracy (\%) & Macro F1 \\
\hline
CDL history & 19 & 80.59 & 0.8080 \\
CDL + NDVI & 34 & \textbf{82.33} & \textbf{0.8245} \\
CDL + SMAP & 23 & 80.66 & 0.8086 \\
CDL + NDVI + SMAP & 38 & 82.26 & 0.8239 \\
\hline
\end{tabular}
\caption{Feature-group comparison on the balanced 2022 validation sample
($n=500{,}000$). All configurations used the same training and validation
years. These are validation results, not comparisons on the 2023 test set.}
\label{tab:ablation}
\end{table}

NDVI produced the largest improvement over crop history
(Table~\ref{tab:ablation}): accuracy increased by 1.74 percentage points.
Winter wheat showed the largest class-specific gain, with F1 rising from
0.849 to 0.881. This is consistent with the earlier seasonal peak seen in
the separate phenology analysis, although the comparison does not isolate
which NDVI feature caused the gain. Adding SMAP to CDL alone increased
accuracy by 0.07 percentage points. The full model scored slightly below
CDL plus NDVI. These observed differences support NDVI's value in this
validation year; they do not establish statistical significance or the same
ordering in other years.

\subsection{Test-year errors and crop-history patterns}

The full model achieved 79.21\% accuracy and 0.7914 macro F1 in 2023
(Figure~\ref{fig:prediction}). F1 was 0.893 for Other, 0.813 for
winter wheat, 0.726 for corn, and 0.735 for soybean. The largest individual
confusion was soybean classified as corn: 30,935 of the 125,000 soybean
samples, compared with 9,894 corn samples classified as soybean. The model
therefore retained substantial confusion between the two summer crops
despite the additional seasonal observations.

To examine performance across crop histories, the feature builder assigned
regimes using each target year's preceding historical window. These are
separate from the 2015--2024 rotation map. A regular history required an
alternation score of at least 0.7 and a distance of at most three from a
canonical corn--soybean sequence. Monoculture took precedence when the
longest repeated-crop run reached seven years or the corn or soybean
fraction reached 0.8; remaining histories were labeled irregular.

Accuracy was 95.55\% for monoculture histories ($n=124{,}318$), 87.41\% for
regular rotation ($n=65{,}574$), and 70.92\% for irregular histories
($n=310{,}108$). Different class mixtures across these groups prevent a
causal interpretation of the accuracy gaps. The regular group's high accuracy
also masks uneven performance across classes: corn and soybean F1 scores were
0.879 and 0.907, while four-class macro F1 was 0.559.

\subsection{Interpreting the fitted model}

To examine which inputs influenced the fitted model, we applied SHAP
TreeExplainer \citep{lundberg2017shap} to a random sample of 1,000 test
observations. Mean absolute attribution ranked the previous
year's CDL code first, followed by mid-season NDVI, NDVI peak timing, and
late- and early-season NDVI means. Growing-season mean soil moisture ranked
eighth. These rankings describe the fitted model's use of its inputs.
